% ==============================================================================
% Plantilla Institucional de Trabajo Terminal — UPIIZ IPN
% Archivo: capitulos/protocolo/08-productos-resultados-esperados.tex
% ==============================================================================
% ESTRUCTURA NORMATIVA DE PROTOCOLO: CAPÍTULO 8 - PRODUCTOS O RESULTADOS ESPERADOS
% ==============================================================================
% LÍMITE INSTITUCIONAL: Máximo 1 página.
%
% Requisitos normativos (Anexo 1):
% 1. Productos o resultados esperados.
% 2. Documentos que se generarán.
% 3. Pruebas o experimentos que se prevé realizar para validar los resultados de la implementación.
% 4. Correspondencia con la Línea de Trabajo seleccionada.
%
% Criterio de redacción:
% - No forzar divisiones artificiales como "tangibles/intangibles".
% ==============================================================================

\chapter{Productos o resultados esperados}
\label{cap:proto-productos-esperados}

% [INSTRUCCIÓN PARA EL ALUMNO]:
% Desarrolle este capítulo en un máximo de una página cubriendo los puntos del Anexo 1:

\section{Productos o resultados esperados y correspondencia con la Línea de Trabajo}

% Conforme al Anexo 1, el entregable formal debe corresponder a la Línea de Trabajo seleccionada:
% Línea: \NumeroLineaTrabajo\ - \NombreLineaTrabajo
% Entregable oficial: \EntregableLineaTrabajo

[Indique con claridad los productos o resultados concretos que se obtendrán al término del Trabajo Terminal, precisando su correspondencia con el entregable institucional de la Línea de Trabajo seleccionada (\NumeroLineaTrabajo: \EntregableLineaTrabajo)].

\section{Documentos que se generarán}

[Especifique los documentos formales que se elaborarán como resultado del proyecto, tales como el reporte técnico de Trabajo Terminal I, el reporte final de Trabajo Terminal II, planos mecánicos, diagramas esquemáticos de circuitos, código fuente documentado, manuales de operación o memorias de cálculo].

\section{Pruebas o experimentos previstos para validación}

[Describa las pruebas o experimentos específicos que se prevé realizar para validar cuantitativamente los resultados obtenidos tras la implementación del sistema].
